\section{Descripción de Participantes y Usuarios}
Para diseñar un sistema efectivo es necesario identificar a todas las partes
afectadas por el proyecto y comprender sus necesidades reales. Este capítulo
distingue entre \textbf{participantes} (personas u organizaciones afectadas por el
proyecto pero que no son usuarios directos del sistema) y \textbf{usuarios} (personas
que interactúan directamente con la aplicación), describe sus perfiles, presenta
una visión general del producto desde la perspectiva de cada uno, detalla el
entorno de despliegue, recoge las suposiciones y dependencias del proyecto y
aborda los aspectos económicos.

\subsection{Resumen de los Participantes}
A continuación se enumeran las partes que se ven afectadas por el proyecto sin
ser necesariamente usuarios directos de la aplicación.

\vspace{5mm}

\begin{longtable}[]{@{}
    % Metadatos de la tabla
        p{0.28\columnwidth}
        p{0.22\columnwidth}
        p{0.50\columnwidth}@{}}

    \caption{Resumen de Participantes} \\

    % Configuración de la tabla para repetir encabezados en cada página
    \toprule
    \textbf{Participante} & \textbf{Naturaleza} & \textbf{Implicación en el proyecto} \\
    \midrule
    \endfirsthead % chktex 1

    \toprule
    \textbf{Participante} & \textbf{Naturaleza} & \textbf{Implicación en el proyecto} \\
    \midrule
    \endhead % chktex 1

    \midrule
    \multicolumn{2}{r}{\small\emph{Continúa en la siguiente página}} \\
    \endfoot % chktex 1

    \bottomrule
    \endlastfoot % chktex 1

    % Cuerpo de la tabla
    \textbf{Coordinación de la Casa de la Esperanza} & Cliente / Usuario &
    Promotor del proyecto, principal fuente de requisitos y usuario
    administrador del sistema (Ezequiel). \\
    \midrule

    \textbf{Directiva de la asociación} & Cliente / Usuario & Toma
    decisiones estratégicas sobre la asociación; consume informes y datos
    agregados del sistema. Incluye contabilidad (Miguel) y otros
    responsables de área. \\
    \midrule

    \textbf{Voluntariado operativo} & Usuario & Conjunto de voluntarios que
    cubren turnos en el centro: café, ropa, duchas, lavandería, recepción.
    Heterogéneo en cuanto a alfabetización digital. \\
    \midrule

    \textbf{Voluntariado profesional} & Usuario & Profesionales (abogados,
    médicos, psicólogos, podólogos, asistentes en extranjería) que prestan
    servicio especializado de forma puntual o regular. \\
    \midrule

    \textbf{Amigos atendidos} & Beneficiario / Sujeto de datos & Personas en
    situación de vulnerabilidad atendidas por la asociación. No son usuarios
    del sistema, pero sus datos personales constituyen el núcleo de la
    información gestionada. \\
    \midrule

    \textbf{Gestoría fiscal externa} & Tercero receptor & Recibe la
    información contable que la asociación le remite para la presentación de
    obligaciones fiscales. Interactúa con el sistema de forma indirecta a
    través de exportaciones. \\
    \midrule

    \textbf{Entidades financiadoras} & Tercero receptor & Reciben informes
    de actividad y de gasto con destino a la justificación de subvenciones.
    Incluyen La Caixa, Ayuntamiento de Málaga y otras entidades públicas y
    privadas. \\
    \midrule

    \textbf{Autoridades sanitarias} & Tercero receptor & Receptoras
    eventuales del cuaderno de trazabilidad alimentaria en caso de
    inspección, conforme a la Ley 1/2025. \\
    \midrule
    
    \textbf{Entidades donantes con convenio} & Tercero & Bancosol, Mercadona
    y otras entidades que donan alimentos o bienes a la asociación bajo
    acuerdos formales. \\
    \midrule

    \textbf{Donantes particulares} & Tercero & Personas que
    realizan donaciones puntuales (alimentos, ropa, dinero) sin convenio
    formal. \\
    \midrule

    \textbf{Equipo de desarrollo} & Proveedor & Responsable del desarrollo,
    despliegue y mantenimiento del sistema. Compuesto por un único
    desarrollador en la fase actual del proyecto. \\
    \midrule
    
    \textbf{Render (proveedor de infraestructura)} & Proveedor técnico &
    Plataforma de alojamiento gestionado en la nube donde se despliega la
    aplicación. \\
\end{longtable}

\newpage

\subsection{Perfiles de los Participantes}
A continuación se detallan los perfiles de los participantes que no son
usuarios directos del sistema pero cuyas necesidades, restricciones o
expectativas afectan al diseño.

\subsubsection{Coordinación}
\begin{itemize}
    \item \textbf{Representante:} Ezequiel.
    
    \item \textbf{Responsabilidades:} Dirección operativa diaria de la Casa de la Esperanza,
    relación con voluntariado, decisión última sobre admisión y trato a los
    amigos, autorización de compras significativas, relación con donantes y
    entidades externas.
    
    \item \textbf{Necesidades respecto al sistema:} Visibilidad completa sobre toda la
    operación, trazabilidad por persona y
    por universo, capacidad de delegar acceso granular a otros voluntarios
    manteniendo el control sobre datos sensibles, herramientas que reduzcan la
    carga manual y permitan rendir cuentas a la directiva y a financiadores con
    datos exactos en lugar de porcentajes aproximados
\end{itemize}


\subsubsection{Directiva (Contabilidad)}
\begin{itemize}
    \item \textbf{Representante:} Miguel 
    
    \item \textbf{Responsabilidades:} Contabilidad de la asociación, relación con la
    gestoría fiscal, validación de gastos, control financiero, informes a
    órganos de gobierno y financiadores.
    
    \item \textbf{Necesidades respecto al sistema:} Importación automatizada de extractos
    bancarios, mapeo de conceptos a cuentas contables, generación de asientos
    sin trabajo manual repetitivo, acceso desde el domicilio a los datos
    contables, capacidad de consulta autónoma sin depender de la gestoría,
    trazabilidad de gasto por persona atendida.
\end{itemize}


\subsubsection{Gestoría Fiscal Externa}
\begin{itemize}
    \item \textbf{Responsabilidades:} Cierre fiscal de la asociación.
    
    \item \textbf{Necesidades respecto al sistema:}  Recibir información estructurada,
    exportable y trazable de la actividad económica de la asociación.
    
    \item \textbf{Restricciones:} El sistema debe poder exportar a
    formatos estándar (CSV, PDF) compatibles con el flujo de la gestoría. La
    gestoría no es un usuario del sistema y no se contempla integración
    técnica en esta versión.
\end{itemize}


\subsubsection{Autoridades Sanitarias}
\begin{itemize}
    \item \textbf{Necesidades respecto al sistema:} En caso de inspección, recibir el
    cuaderno de entradas y salidas alimentarias con la trazabilidad completa
    exigida por la Ley 1/2025 (lotes, caducidades, temperaturas, prioridad
    legal, convenios).
    
    \item \textbf{Implicaciones de diseño:} El cuaderno debe ser exportable en formato
    imprimible y debe garantizar la inmutabilidad de los registros una vez
    cerrado el periodo correspondiente.
\end{itemize}


\subsubsection{Equipo de Desarrollo}
\begin{itemize}
    \item \textbf{Composición actual:} Un único desarrollador.
    
    \item \textbf{Implicaciones de diseño:} La arquitectura debe favorecer la
    mantenibilidad por una sola persona, la simplicidad operativa, el uso de
    herramientas y librerías ampliamente documentadas y la cobertura de
    pruebas suficiente para detectar regresiones sin necesidad de equipos de
    QA dedicados.
\end{itemize}


\newpage

\subsection{Perfiles de los Usuarios}
Los usuarios directos del sistema se agrupan en cinco perfiles, cada uno con
un rol asociado en el modelo de control de acceso definido en RF2.1. La
relación entre perfil de usuario y rol no es uno a uno: un mismo voluntario
puede acumular varios roles si su nivel de implicación lo justifica.

\subsubsection{Adminitrador}
\begin{itemize}
    \item \textbf{Rol del sistema:} \texttt{ADMIN}.
    
    \item \textbf{Quién asume este perfil:} Coordinación (Ezequiel) y,
    personas de máxima confianza de la directiva (por ejemplo, Miguel).
    
    \item \textbf{Permisos:}
    \begin{itemize}
        \item Crear, editar y desactivar cuentas de voluntario.
        \item Asignar y revocar roles, niveles de seguridad y funciones específicas.
        \item Habilitar accesos temporales a voluntarios profesionales sobre amigos concretos.
        \item Acceso completo a todos los datos del sistema.
        \item Acceso a la gestión financiera y creación de asientos contables.
        \item Exportación de informes para entidades financiadoras y autoridades.
        \item Acceso al registro de auditoría.
    \end{itemize}
\end{itemize}


\subsubsection{Voluntario Operativo}
\begin{itemize}
    \item \textbf{Rol del sistema:} \texttt{OPERATIVO}.
    
    \item \textbf{Quién asume este perfil:} El voluntariado del centro;
    personas que cubren turnos de mañana o tarde en tareas como recepción,
    café, ropa, duchas, lavandería y registro de servicios.
    
    \item \textbf{Permisos:}
    \begin{itemize}
        \item Consultar y registrar fichas de amigo en niveles de confianza: \texttt{BASICO} (sin
        acceso a datos sensibles del nivel superior salvo que su nivel de seguridad lo
        habilite explícitamente) o \texttt{PRIVILEGIADO}.
        \item Agendar y registrar la prestación de servicios.
        \item Realizar entregas de ropa y operar el flujo de lavandería.
        \item Registrar entradas y salidas alimentarias.
        \item Consultar el stock de armario.
        \item Sin acceso por defecto al módulo contable ni a notas privadas.
    \end{itemize}
\end{itemize}


\subsubsection{Voluntario Profesional}
\begin{itemize}
    \item \textbf{Rol del sistema:} \texttt{PROFESIONAL}.
    
    \item \textbf{Quién asume este perfil:} Profesionales que prestan servicio
    especializado de forma puntual o regular: abogados, médicos, psicólogos,
    podólogos, asistentes en extranjería.
    
    \item \textbf{Permisos:}
    \begin{itemize}
        \item Consultar fichas de los amigos a los que han sido habilitados por la
        coordinación, incluida la información sensible del área de su
        especialidad (notas jurídicas para el abogado, notas médicas para el
        personal sanitario, etc.).
        \item Registrar nuevas notas y observaciones en el ámbito de su especialidad.
        \item Sin acceso a información de áreas distintas de la suya, salvo
        habilitación explícita y temporal por parte del administrador.
    \end{itemize}
\end{itemize}


\subsubsection{Voluntario Puente}
\begin{itemize}
    \item \textbf{Rol del sistema:} \texttt{PUENTE}.
    
    \item \textbf{Quién asume este perfil:} Personas que se encuentran en una posición
    intermedia entre amigo y voluntario. Caso de referencia: Mohamed.
    
    \item \textbf{Permisos:}
    \begin{itemize}
        \item Permisos básicos de voluntario operativo limitados a las tareas
        específicas que se le asignen (típicamente: consulta de stock,
        registro de su propia ropa en lavandería, agendamiento de duchas).
        \item Sin acceso a datos sensibles de otros amigos ni al módulo contable.
        \item No puede consultar datos sensibles de
        sí mismo a través de su sesión de voluntario, sino que debe seguir el
        procedimiento estándar previsto para los amigos.
    \end{itemize}
\end{itemize}


\subsubsection{Agente Externo Temporal}
\begin{itemize}
    \item \textbf{Rol del sistema:} \texttt{AGENTE\_TEMP}.
    
    \item \textbf{Quién asume este perfil:} Profesionales o colaboradores externos a
    quienes la coordinación necesita conceder acceso muy puntual a una
    información concreta, durante una ventana de tiempo limitada y sin
    incorporarlos a la base estable de voluntarios. Por ejemplo, un abogado
    externo invitado a revisar un expediente jurídico específico.
    
    \item \textbf{Permisos:}
    \begin{itemize}
        \item Acceso de solo lectura a las entidades específicas habilitadas por el administrador.
        \item La habilitación incluye una fecha de caducidad explícita pasada la
        cual el acceso queda revocado automáticamente.
        \item Toda actividad del agente queda registrada en el log de auditoría con
        marcado especial.
    \end{itemize}
\end{itemize}


\subsubsection{Tabla resumen de roles y permisos}

\begin{longtable}[]{
    @{}p{0.60\columnwidth}
    p{0.05\columnwidth}
    p{0.05\columnwidth}
    p{0.05\columnwidth}
    p{0.05\columnwidth}
    p{0.06\columnwidth}@{}}

    \caption{Resumen de Roles y Permisos} \\

    % Configuración de la tabla para repetir encabezados en cada página
    \toprule
    \textbf{Operación} & \textbf{Ad} & \textbf{Op} & \textbf{Pro} & \textbf{Pu} & \textbf{Atemp.} \\
    \midrule
    \endfirsthead % chktex 1

    \toprule
    \textbf{Operación} & \textbf{Ad} & \textbf{Op} & \textbf{Pro} & \textbf{Pu} & \textbf{Atemp.} \\
    \midrule
    \endhead % chktex 1

    \midrule
    \multicolumn{2}{r}{\small\emph{Continúa en la siguiente página}} \\
    \endfoot % chktex 1

    \bottomrule
    \endlastfoot % chktex 1

    % Cuerpo de la tabla
    AGENDAR\_SERVICIOS & $\bullet$ & $\bullet$ & $\bullet$ & $\circ$  & --- \\
    CONCEDER\_HABILITACION\_TEMPORAL & $\bullet$ & --- & --- & --- & --- \\
    CONSULTAR\_LOG\_AUDITORIA & $\bullet$ & --- & --- & --- & --- \\
    CONSULTAR\_NOTAS\_AMIGO & $\bullet$ & $\circ$  & $\circ$  & --- & $\circ$  \\
    CREAR\_ASIENTOS & $\bullet$ & $\circ$ & --- & --- & --- \\
    CREAR\_NOTA\_PRIVADA & $\bullet$ & $\bullet$  & $\bullet$  & $\bullet$ & $\bullet$  \\
    ENTREGAR\_ROPA & $\bullet$ & $\bullet$ & --- & --- & --- \\
    ESCALAR\_NIVEL\_3 & $\bullet$ & $\circ$  & --- & --- & --- \\
    EXPORTAR\_INFORMES & $\bullet$ & $\circ$  & --- & --- & $\circ$  \\
    GESTIONAR\_ARMARIO & $\bullet$ & $\bullet$ & $\bullet$ & $\bullet$ & $\circ$  \\
    GESTIONAR\_CUENTAS\_USUARIO & $\bullet$ & --- & --- & --- & --- \\
    IMPORTAR\_DATOS & $\bullet$ & $\bullet$ & --- & --- & --- \\
    IMPORTAR\_EXTRACTOS & $\bullet$ & --- & --- & --- & --- \\
    OPERAR\_LAVANDERIA & $\bullet$ & $\bullet$ & --- & $\circ$  & --- \\
    REGISTRAR\_ENTREGA\_ALIMENTARIA & $\bullet$ & $\bullet$ & --- & --- & --- \\
    REGISTRAR\_GASTO\_AMIGO & $\bullet$ & $\circ$  & --- & --- & --- \\
    REGISTRAR\_SALIDA\_ALIMENTARIA & $\bullet$ & $\bullet$ & --- & --- & --- \\
    REGISTRAR\_AMIGO & $\bullet$ & $\bullet$ & --- & --- & --- \\
\end{longtable}

\textbf{Leyenda:} 
\par
$\bullet$  permiso completo; 
\par
$\circ$  permiso condicional según
asignación explícita del administrador; 
\par
--- sin permiso.


